% 星际行星（综述）
% license CCBYSA3
% type Wiki

本文根据 CC-BY-SA 协议转载翻译自维基百科\href{https://en.wikipedia.org/wiki/Rogue_planet}{相关文章}。

流浪行星（rogue planet），亦称自由漂浮行星（free-floating planet，FFP）或孤立的行星质量天体（isolated planetary-mass object，iPMO），是指一种具有行星质量、但在引力上不与任何恒星或棕矮星发生束缚关系的星际天体$^{[1][2][3][4]}$。

流浪行星可能起源于其形成所在的行星系统，并在随后被抛射出去；它们也可能在行星系统之外独立形成。仅银河系中就可能存在数十亿到数万亿颗流浪行星，这一数量范围有望由即将投入使用的南希·格雷斯·罗曼空间望远镜进一步加以精确限定$^{[5][6]}$。流浪行星进入太阳系的概率极低，更不用说对地球生命构成直接威胁——在未来一千年内发生这种情况的概率约为一万亿分之一$^{[7]}$。

一些行星质量天体可能以类似恒星的方式形成，国际天文学联合会已提议将此类天体称为亚棕矮星（sub-brown dwarfs）$^{[8]}$。一个可能的例子是 Cha~110913$-$773444，它可能是在被抛射后成为一颗流浪行星，也可能是独立形成并成为一颗亚棕矮星$^{[9]}$。
\subsection{术语}
最早的两篇发现论文分别使用了“孤立的行星质量天体”（iPMO）$^{[10]}$和“自由漂浮行星”（FFP）$^{[11]}$这两个名称。大多数天文学论文使用其中之一$^{[12][13][14]}$。“流浪行星”这一术语更常用于引力微透镜研究，而这类研究也经常使用 FFP 这一术语$^{[15][16]}$。面向公众的新闻稿则可能采用不同的名称。例如，2021 年发现至少 70 颗 FFP 时，不同新闻稿分别使用了“流浪行星”$^{[17]}$、“无恒星行星”$^{[18]}$、“游荡行星”$^{[19]}$以及“自由漂浮行星”$^{[20]}$等称呼。
\subsection{发现}
孤立的行星质量天体（isolated planetary-mass objects，iPMO）最早于 2000 年由英国研究团队 Lucas \& Roche 使用英国红外望远镜（UKIRT）在猎户座星云中发现$^{[11]}$。同年，西班牙研究团队 Zapatero Osorio 等人利用凯克望远镜的光谱观测，在 $\sigma$~猎户座星团中发现了 iPMO$^{[10]}$。对猎户座星云中这些天体的光谱研究于 2001 年发表$^{[21]}$。目前，这两个欧洲研究团队均因其准同步发现而受到认可$^{[22]}$。1999 年，日本研究团队 Oasa 等人在变色龙座 I 云中发现了一些天体$^{[23]}$，这些天体随后在 2004 年由美国研究团队 Luhman 等人通过光谱观测得到了确认$^{[24]}$。
\subsection{观测}
\begin{figure}[ht]
\centering
\includegraphics[width=6cm]{./figures/35a477b3258b7d43.png}
\caption{2021 年在上天蝎座与蛇夫座之间的区域发现了 115 个潜在的流浪行星} \label{fig_RoPl_1}
\end{figure}
发现自由漂浮行星主要有两种技术：直接成像和引力微透镜。
\subsubsection{引力微透镜}
日本大阪大学的天体物理学家角广隆宏（Takahiro Sumi）及其同事，隶属于天体物理学引力微透镜观测计划（Microlensing Observations in Astrophysics）和光学引力透镜实验（Optical Gravitational Lensing Experiment）合作组，于 2011 年发表了他们关于引力微透镜的研究成果。他们使用位于新西兰约翰山天文台的 1.8 米（5 英尺 11 英寸）MOA-II 望远镜，以及位于智利拉斯坎帕纳斯天文台的华沙大学 1.3 米（4 英尺 3 英寸）望远镜，对银河系中约 5000 万颗恒星进行了观测。他们共发现了 474 次微透镜事件，其中有 10 次事件的持续时间足够短，表明它们可能是质量约为木星质量、且在其附近没有伴随恒星的行星。研究人员根据这些观测结果估计，在银河系中，每一颗恒星大约对应近两颗木星质量的流浪行星$^{[25][26][27]}$。另有一项研究提出了一个大得多的数量估计，认为银河系中的流浪行星数量可能是恒星数量的多达 100{,}000 倍，不过该研究所涵盖的假想天体质量远小于木星$^{[28]}$。2017 年，华沙大学天文台的 Przemek~Mróz 及其同事在统计样本数量为 2011 年研究六倍的基础上开展的一项研究表明，银河系中每一颗主序星所对应的木星质量自由漂浮行星或宽轨道行星的上限为 0.25 颗$^{[29]}$。

2020 年 9 月，天文学家利用引力微透镜技术首次报告探测到一颗地球质量的流浪行星（命名为 OGLE-2016-BLG-1928），该天体不与任何恒星发生引力束缚，并在银河系中自由漂浮$^{[16][30][31]}$。
\subsubsection{直接成像}
\begin{figure}[ht]
\centering
\includegraphics[width=6cm]{./figures/250326c529d9a895.png}
\caption{使用斯皮策空间望远镜观测到的低温行星质量天体 WISE~J0830+2837（以橙色天体标示）。其温度约为 300--350~K（27--77\,$^\circ$C；80--170\,$^\circ$F）。} \label{fig_RoPl_2}
\end{figure}
微透镜行星只能通过微透镜事件本身进行研究，这使得对行星性质的刻画变得困难。因此，天文学家转而研究通过直接成像方法发现的孤立行星质量天体（iPMO）。例如，要确定一颗棕矮星或 iPMO 的质量，需要其光度和年龄等信息$^{[32]}$。然而，确定低质量天体的年龄已被证明是困难的。因此，绝大多数 iPMO 被发现于附近年轻的恒星形成区中，这些区域的年龄为天文学家所知。这些天体的年龄小于 200~Myr，质量较大（$>5\,M_J$）$^{[4]}$，并且属于 $L$ 矮星和 $T$ 矮星$^{[33][34]}$。然而，也存在一个规模虽小但正在增长的低温且年老的 $Y$ 矮星样本，其估计质量为 $8$--$20\,M_J$ $^{[35]}$。光谱型为 $Y$ 的附近流浪行星候选体包括距离为 $7.27\pm0.13$ 光年的 WISE~0855$-$0714$^{[36]}$。如果能够通过更精确的测量对这一 Y 矮星样本进行刻画，或找到更好地确定其年龄的方法，那么年老且低温的 iPMO 数量很可能会显著增加。

最早的 iPMO 于 21 世纪初通过对年轻恒星形成区的直接成像而被发现$^{[37][10][21]}$。这些通过直接成像发现的 iPMO 很可能以类似恒星的方式形成（有时被称为亚棕矮星）。也可能存在以行星方式形成、随后被抛射出去的 iPMO。然而，这类天体在运动学上应当不同于其诞生的恒星形成区，不应被环星盘所包围，并且应具有较高的金属丰度$^{[22]}$。在年轻恒星形成区中发现的 iPMO 中，没有任何一个相对于其恒星形成区表现出较高的速度。对于年老的 iPMO，低温天体 WISE~J0830+2837$^{[38]}$ 的切向速度 $V_{\mathrm{tan}}$ 约为 $100\,\mathrm{km/s}$，这一数值虽高，但仍与其在银河系中形成的情形相一致。对于 WISE~1534$-$1043$^{[39]}$，由于其约为 $200\,\mathrm{km/s}$ 的高 $V_{\mathrm{tan}}$，一种替代情景将其解释为被抛射出的系外行星，但其颜色却表明它是一颗年老、低金属丰度的棕矮星。大多数研究大质量 iPMO 的天文学家认为，它们代表了恒星形成过程的低质量端$^{[22]}$。

天文学家利用赫歇尔空间天文台和甚大望远镜观测了一颗非常年轻的自由漂浮行星质量天体 OTS~44，并证明了典型的类恒星形成模式所具有的过程同样适用于质量低至数个木星质量的孤立天体。赫歇尔的远红外观测显示，OTS~44 被一个质量至少为 10 个地球质量的盘所包围，因此最终可能形成一个大型卫星系统$^{[40]}$。利用甚大望远镜上的 SINFONI 光谱仪对 OTS~44 进行的光谱观测揭示，该盘正在积极地吸积物质，这与年轻恒星周围的盘相似$^{[40]}$。
\subsubsection{双星}
\begin{figure}[ht]
\centering
\includegraphics[width=6cm]{./figures/2fb0b35b79287356.png}
\caption{2MASS~J1119–1137AB，即首个被发现的行星质量双星系统，位于 TW~Hydrae 协会中} \label{fig_RoPl_3}
\end{figure}
\begin{figure}[ht]
\centering
\includegraphics[width=6cm]{./figures/4a9c799a9b362b83.png}
\caption{JuMBO~29，一个候选的 $12.5+3\,M_J$ 行星质量双星系统，分离距离为 135~AU，位于猎户座星云中} \label{fig_RoPl_4}
\end{figure}